\begin{abstract}
A recent study reported $0.9047$ balanced accuracy (BA) for
fine-tuned LaBraM on the Komarov--Ko--Jung resting-state stress
cohort \citep{wang2025stress,komarov2020stress}, suggesting
near-clinical-grade EEG-foundation-model (FM) performance on
cross-subject decoding. We ask when such claims should be expected
to hold. Across three FMs (LaBraM, CBraMod, REVE) and four datasets
(UCSD Stress, EEGMAT, ADFTD, TDBRAIN) we show that (i) frozen FM
representations encode subject identity $10$--$50\times$ more
strongly than task labels in all 12 model-dataset cells; (ii)
trial-to-subject-level cross-validation contributes a $1$--$15$\,pp
component of the reduction to our reproduced $0.52$--$0.58$ Stress
band, with the remaining gap to $0.9047$ attributable to unreported
preprocessing or hyperparameter variations in the reference pipeline; (iii) seven
architectures spanning five orders of magnitude in parameter count
converge to a single $5$\,pp window on Stress; and (iv) a controlled
paired within-subject comparison --- EEGMAT rest-vs-arithmetic
(anchored by Klimesch alpha desynchronisation) versus UCSD Stress
longitudinal DSS (no published within-subject EEG correlate and no
cluster-corrected dataset-level spectral contrast in our
pre-registered anchor test) --- produces opposite rescue outcomes
(BA $0.73$ vs $0.17$--$0.43$), mirrored by cross-model neural-band
consensus. We call this pattern the \emph{Subject-Dominance Limit}
(SDL): cross-subject FM classification is bounded by the
within-subject neural contrast strength of the target label, not by
model capacity. We release a pre-benchmark contrast-anchoring
protocol.
\end{abstract}
